\section{Experiments}
\subsection{Datasets and Few-Shot Protocol}
We evaluate on MVTec AD and VisA. For each class, $k\in\{1,2,4,8\}$ normal support images are sampled with fixed seeds where applicable. The main full-scale conformal experiment uses VisA with 12 classes, $k\in\{4,8\}$, five seeds, and four corruptions: Gaussian noise, blur, brightness/contrast shift, and JPEG compression. This yields 480 cases and 56,000 image-level rows after capping evaluation to at most 120 images per case.

\subsection{Baselines and Variants}
We compare against PatchCore/AnomalyDINO-style memory-bank baselines, a SubspaceAD-style PCA residual baseline, scalar Platt, Vector Platt, Shift-Aware Platt, weighted Platt, LOIO conformal, and weighted conformal. Official SubspaceAD representative runs are used as a novelty guardrail: they show that frozen DINOv2 PCA/subspace ranking is already strong and should not be claimed as our novelty.

\subsection{Metrics}
Metrics are separated by role. Ranking quality is measured by AUROC and AP using the raw PCA/subspace score. Pixel-level localization is reported by pixel AUROC and AU-PRO when masks are available. Reliability is measured by expected calibration error (ECE), Brier score, and negative log-likelihood (NLL) using the selected probability-like reliability view. Efficiency is reported as reference/model storage and cached-feature latency. Robustness diagnostics include corruption shift and label-aware PCA-residual FGSM surrogate tests.

\subsection{Claim Discipline}
The paper uses conservative claim rules. A method that changes only reliability probabilities should not be credited with ranking improvement. Adversarial results are reported as fragility diagnostics. Conformal and gated mechanisms are positioned as reliability-routing components rather than as first-in-field conformal or gated AD claims.
